\chapter{Brain and Nervous System Diseases}
\label{ch:brain}
The brain is the command center of the human body, weighing approximately 1.4 kg and containing roughly 86 billion neurons. It controls voluntary and involuntary actions, processes sensory information, governs cognition and emotion, and maintains homeostasis. The brain is protected by the skull, three meningeal layers, and cerebrospinal fluid. Diseases of the brain and nervous system represent some of the most debilitating conditions in medicine, affecting movement, sensation, cognition, and vital functions.

%-------------------------------------------------------------
\section{Brain Tumors}
%-------------------------------------------------------------

%-------------------------------------------------------------

%-------------------------------------------------------------

%-------------------------------------------------------------

Brain tumors represent a heterogeneous group of primary and metastatic central nervous system neoplasms classified by cell of origin, histology, and WHO grade. Primary brain tumors originate within cranial parenchymal cells, meninges, or cranial nerves, whereas metastatic lesions arise from systemic malignancies spreading hematogenously to the neuroaxis. Progressive tumor expansion increases intracranial pressure, producing focal neurological deficits, cognitive decline, and seizures requiring multimodality therapy.

%-------------------------------------------------------------

%-------------------------------------------------------------

%-------------------------------------------------------------

%-------------------------------------------------------------

\label{sec:Brain Tumors}
%-------------------------------------------------------------%-------------------------------------------------------------

\subsection{Glioblastoma}
\index{Glioblastoma}
\label{sec:Glioblastoma}

Glioblastoma is a WHO grade 4 astrocytoma and the most common and aggressive primary malignant brain tumor in adults. The condition results from rapid infiltrative growth with microvascular proliferation and areas of necrosis; the IDH-wildtype molecular profile is characteristic. Clinical features depend on tumor location but typically include progressive headache, seizures, focal neurological deficits, and cognitive changes. Diagnosis is based on MRI showing a ring-enhancing lesion with surrounding edema and mass effect. Management follows the Stupp protocol: maximal safe surgical resection followed by concurrent temozolomide chemotherapy and radiotherapy, then adjuvant temozolomide; MGMT promoter methylation status predicts response to temozolomide; tumor-treating fields (TTFields) have shown incremental survival benefit. Prognosis is poor despite aggressive treatment, with median survival of approximately 15 months.

\subsection{Medulloblastoma}
\index{Medulloblastoma}
\label{sec:Medulloblastoma}
Medulloblastoma is the most common malignant brain tumor of childhood, arising in the cerebellum or posterior fossa, and belongs to the group of embryonal tumors classified into molecular subgroups (WNT, SHH, Group 3, Group 4) with distinct prognoses. The condition results from malignant transformation of primitive neuroectodermal cells. Clinical features include headache, vomiting (especially morning), ataxia, and signs of hydrocephalus from fourth ventricle obstruction. Diagnosis involves MRI showing a posterior fossa mass and cerebrospinal fluid cytology for staging. Management includes maximal surgical resection, craniospinal irradiation, and adjuvant chemotherapy; risk-adapted therapy aims to reduce long-term neurocognitive and endocrine sequelae in survivors. Prognosis varies by molecular subgroup, with the WNT subgroup having the best prognosis.

\subsection{Meningioma}
\index{Meningioma}
\label{sec:Meningioma}
Meningioma is the most common primary intracranial tumor, arising from arachnoid cap cells; it is usually benign (WHO grade 1), slow-growing, and attached to the dura; meningiomas are more common in women and associated with prior radiation, neurofibromatosis type 2, and high progesterone receptor expression. The condition results from neoplastic proliferation of arachnoid cap cells. Most are asymptomatic and discovered incidentally on imaging; symptomatic meningiomas cause headaches, seizures, or focal deficits depending on location. Diagnosis is based on MRI revealing a well-circumscribed, homogeneously enhancing dural-based mass with a dural tail. Management involves observation for small, asymptomatic lesions; surgical resection is the mainstay for symptomatic or growing tumors; atypical (grade 2) and anaplastic (grade 3) meningiomas require adjuvant radiotherapy. Prognosis is excellent for benign meningiomas.

\subsection{Neurofibromatosis Type 1}
\index{Neurofibromatosis Type 1}
\label{sec:Neurofibromatosis Type 1}
Neurofibromatosis type 1 (NF1, von Recklinghausen disease) is an autosomal dominant neurocutaneous disorder caused by mutations in the NF1 gene encoding neurofibromin, a Ras-GTPase tumor suppressor, affecting 1 in 3,000 births. The condition results from loss of neurofibromin function, leading to uncontrolled Ras signaling and increased cell proliferation. Clinical features include caf\'{e}-au-lait spots ($\ge$ 6 spots $>$ 5 mm in children, $>$ 15 mm in adults), axillary and inguinal freckling, Lisch nodules (iris hamartomas), neurofibromas (cutaneous, subcutaneous, plexiform), optic pathway gliomas, osseous dysplasia, and increased risk of malignant peripheral nerve sheath tumors. Diagnosis is based on NIH consensus criteria ($\ge$ 2 of 7 features). Management involves annual ophthalmologic and neurologic monitoring, surgical excision of symptomatic neurofibromas, MEK inhibitor therapy (selumetinib) for symptomatic plexiform neurofibromas, and surveillance for malignancies. Prognosis is variable; life expectancy is reduced by 10--15 years, primarily due to malignancy and cardiovascular disease.

\subsection{Pituitary Adenoma}
\index{Pituitary Adenoma}
\label{sec:Pituitary Adenoma}
Pituitary adenomas are benign neoplasms of the anterior pituitary gland, classified as microadenomas (less than 10 mm) or macroadenomas (10 mm or greater). The condition results from clonal proliferation of pituitary cells. Functioning adenomas secrete excess hormones: prolactinomas (galactorrhea, amenorrhea, hypogonadism), GH-secreting adenomas (acromegaly or gigantism), ACTH-secreting adenomas (Cushing's disease), and TSH-secreting adenomas (central hyperthyroidism); non-functioning adenomas present with mass effects including bitemporal hemianopia from optic chiasm compression, headaches, and hypopituitarism. Clinical features depend on hormone secretion and tumor size. Diagnosis involves hormonal evaluation and MRI of the sella. Management varies: prolactinomas respond to dopamine agonists (cabergoline); other functioning and non-functioning adenomas typically require transsphenoidal surgery, with radiation for residual or recurrent disease. Prognosis is generally favorable.

\subsection{Schwannoma}
\index{Schwannoma}
\label{sec:Schwannoma}
Schwannomas are benign nerve sheath tumors arising from Schwann cells; vestibular schwannomas (acoustic neuromas) are the most common intracranial schwannomas. The condition results from neoplastic proliferation of Schwann cells. Bilateral vestibular schwannomas are pathognomonic of neurofibromatosis type 2. Clinical features of vestibular schwannoma include unilateral sensorineural hearing loss, tinnitus, and imbalance. Diagnosis is based on MRI showing a cerebellopontine angle mass with an ice-cream cone appearance in the internal auditory canal. Management options include observation with serial imaging, stereotactic radiosurgery (Gamma Knife), and microsurgical resection; facial nerve preservation is a key surgical goal. Prognosis is excellent for benign schwannomas.

%-------------------------------------------------------------
\section{Cerebrovascular Diseases}
%-------------------------------------------------------------

%-------------------------------------------------------------

%-------------------------------------------------------------

%-------------------------------------------------------------

%-------------------------------------------------------------

%-------------------------------------------------------------

%-------------------------------------------------------------

%-------------------------------------------------------------

\label{sec:Cerebrovascular Diseases}
%-------------------------------------------------------------%-------------------------------------------------------------

\subsection{Cerebral Aneurysm}
\index{Cerebral Aneurysm}
\label{sec:Cerebral Aneurysm}
A cerebral aneurysm is a focal dilation of a cerebral artery, most commonly at bifurcation sites in the Circle of Willis; saccular (berry) aneurysms are the most prevalent type. The condition results from weakness in the arterial wall. Risk factors include hypertension, smoking, family history, connective tissue disorders (e.g., Ehlers-Danlos syndrome type II, autosomal dominant polycystic kidney disease), and female sex. Most are asymptomatic and discovered incidentally, but rupture causes subarachnoid hemorrhage with high mortality (approximately 40\%). Diagnosis is by CT angiography, MR angiography, or digital subtraction angiography. Management of unruptured aneurysms may involve observation or prophylactic treatment depending on size, location, and patient factors; ruptured aneurysms require urgent surgical clipping or endovascular coiling. Prognosis is good for unruptured aneurysms but poor following rupture.

\subsection{Cerebral Venous Thrombosis}
\index{Cerebral Venous Thrombosis}
\label{sec:Cerebral Venous Thrombosis}
Cerebral venous thrombosis is a thrombosis of the dural venous sinuses or cerebral veins that leads to impaired venous drainage, increased intracranial pressure, and potential venous infarction; it accounts for less than 1\% of all strokes but disproportionately affects young adults and women, particularly those using oral contraceptives or during pregnancy and the puerperium. The condition results from thrombotic occlusion of cerebral veins or dural sinuses. Clinical features vary widely and may include headache, seizures, focal neurological deficits, papilledema, and altered consciousness. Diagnosis requires CT or MR venography showing a filling defect in the venous sinuses. Management consists of anticoagulation (even in the presence of hemorrhagic venous infarction), management of raised intracranial pressure, and treatment of seizures. Prognosis is generally favorable with prompt diagnosis and treatment.

\subsection{Hemorrhagic Stroke}
\index{Hemorrhagic Stroke}
\label{sec:Hemorrhagic Stroke}
Hemorrhagic stroke is a cerebrovascular event resulting from rupture of a blood vessel within the brain, leading to intraparenchymal or subarachnoid hemorrhage. Intracerebral hemorrhage (ICH) is often associated with chronic hypertension, cerebral amyloid angiopathy, vascular malformations, or coagulopathy; subarachnoid hemorrhage (SAH) most commonly results from rupture of a cerebral aneurysm. The condition results from spontaneous rupture of a cerebral blood vessel. Clinical features include sudden severe headache, vomiting, decreased consciousness, and focal neurological deficits that may progress rapidly; SAH presents with thunderclap headache, meningismus, and loss of consciousness. Diagnosis is based on CT head as the initial diagnostic study. Management of ICH focuses on blood pressure control, reversal of anticoagulation, and neurosurgical intervention when indicated; SAH requires urgent neurosurgical clipping or endovascular coiling of the aneurysm with nimodipine to prevent vasospasm. Prognosis depends on hemorrhage location, size, and timeliness of intervention.

\subsection{Ischemic Stroke}
\index{Ischemic Stroke}
\label{sec:Ischemic Stroke}
Ischemic stroke is a focal brain infarction caused by occlusion of a cerebral artery by a thrombus or embolus, accounting for approximately 85\% of all strokes. The condition results from thrombotic or embolic arterial occlusion leading to focal brain ischemia and infarction. Common etiologies include atherosclerotic large-vessel disease, cardioembolism (e.g., atrial fibrillation), small-vessel occlusive disease (lacunar stroke), and cryptogenic causes. Clinical features depend on the vascular territory affected: middle cerebral artery occlusion causes contralateral hemiparesis, hemisensory loss, and aphasia (if dominant hemisphere); posterior circulation strokes produce visual field defects, ataxia, and crossed deficits. Diagnosis involves urgent non-contrast CT to exclude hemorrhage, followed by CT angiography and perfusion imaging. Management includes intravenous thrombolysis with alteplase (within 4.5 hours) and mechanical thrombectomy for large-vessel occlusions (up to 24 hours in selected patients); secondary prevention includes antithrombotic therapy, statins, and risk factor modification. Prognosis depends on infarct size, location, and timeliness of intervention.

\subsection{Transient Ischemic Attack}
\index{Transient Ischemic Attack}
\label{sec:Transient Ischemic Attack}
A transient ischemic attack (TIA) is a temporary episode of neurological dysfunction caused by focal brain, spinal cord, or retinal ischemia without acute infarction; TIAs are critical warning signs, with the risk of subsequent stroke being highest in the first 48 hours (up to 10\%). The condition results from temporary focal ischemia. Clinical features mirror those of ischemic stroke but resolve completely, typically within minutes to hours. Diagnosis involves brain imaging, vascular imaging (carotid ultrasound, CT or MR angiography), cardiac evaluation, and laboratory studies; the ABCD2 score helps stratify short-term stroke risk. Management includes antiplatelet therapy (aspirin, clopidogrel, or combination), statins, antihypertensives, and carotid endarterectomy or stenting for significant ipsilateral carotid stenosis. Prognosis is good with prompt intervention and risk factor modification.

%-------------------------------------------------------------
\section{Demyelinating and White Matter Diseases}
%-------------------------------------------------------------

%-------------------------------------------------------------

%-------------------------------------------------------------

%-------------------------------------------------------------

%-------------------------------------------------------------

%-------------------------------------------------------------

%-------------------------------------------------------------

%-------------------------------------------------------------

\label{sec:Demyelinating and White Matter Diseases}
%-------------------------------------------------------------%-------------------------------------------------------------

\subsection{Central Pontine Myelinolysis}
\index{Central Pontine Myelinolysis}
\label{sec:Central Pontine Myelinolysis}
Central pontine myelinolysis, also known as osmotic demyelination syndrome, is a non-inflammatory demyelinating condition of the pons caused by overly rapid correction of hyponatremia. The condition results from osmotic stress leading to oligodendrocyte damage and demyelination. Clinical features include spastic quadriparesis, pseudobulbar palsy (dysphagia, dysarthria, emotional lability), and in severe cases, locked-in syndrome. Diagnosis is based on MRI showing a characteristic symmetric T2-hyperintense lesion in the central basis pontis. Management is primarily preventive: serum sodium should be corrected at a rate not exceeding 8--10 mmol/L in 24 hours; treatment is supportive once established. Prognosis is variable, with recovery being partial or absent depending on severity.

\subsection{Guillain-Barr\'{e} Syndrome}
\index{Guillain-Barr\'{e} Syndrome}
\label{sec:Guillain-Barré Syndrome}
Guillain-Barr\'{e} syndrome is an acute, immune-mediated polyneuropathy characterized by rapidly progressive, symmetric limb weakness with areflexia; it often follows a respiratory or gastrointestinal infection by 1--4 weeks, with \textit{Campylobacter jejuni}, cytomegalovirus, Epstein-Barr virus, and Zika virus being common triggers. The condition results from an autoimmune attack on peripheral nerves triggered by preceding infection. The most common variant is acute inflammatory demyelinating polyneuropathy (AIDP), while acute motor axonal neuropathy (AMAN) is more prevalent in Asia. Clinical features include rapidly progressive symmetric weakness, areflexia, and possible autonomic dysfunction (tachycardia, blood pressure lability) and respiratory muscle weakness requiring ventilatory support. Diagnosis is based on clinical features and CSF analysis showing albuminocytological dissociation (elevated protein with normal cell count). Management includes intravenous immunoglobulin (IVIg) or plasmapheresis. Prognosis is variable, with most patients recovering; 5--10\% die and 10--20\% have significant residual disability.

%-------------------------------------------------------------
\section{Disorders of Consciousness}
%-------------------------------------------------------------

%-------------------------------------------------------------

%-------------------------------------------------------------

%-------------------------------------------------------------

%-------------------------------------------------------------

%-------------------------------------------------------------

%-------------------------------------------------------------

%-------------------------------------------------------------

\label{sec:Disorders of Consciousness}
%-------------------------------------------------------------%-------------------------------------------------------------

\subsection{Coma}
\index{Coma}
\label{sec:Coma}
Coma is a state of profound unconsciousness in which the patient cannot be aroused and does not respond to external stimuli. The condition results from dysfunction of the ascending reticular activating system in the brainstem or bilateral cerebral hemispheres; common causes include traumatic brain injury, stroke, intracranial hemorrhage, status epilepticus, metabolic encephalopathy, drug intoxication, infections, and hypoxic-ischemic brain injury. Clinical features include unresponsiveness to external stimuli with eyes closed. Diagnosis is based on clinical assessment using the Glasgow Coma Scale (score 3--15), urgent neuroimaging, and laboratory studies. Management targets the underlying cause, with supportive measures including airway management, intracranial pressure monitoring, and prevention of secondary brain injury. Prognosis varies widely depending on the cause and severity of injury.

\subsection{Persistent Vegetative State}
\index{Persistent Vegetative State}
\label{sec:Persistent Vegetative State}
A persistent vegetative state is a condition of wakeful unconsciousness in which the patient has sleep-wake cycles and may open their eyes spontaneously but shows no evidence of awareness of self or environment. The condition results from severe diffuse brain injury, cardiac arrest, or massive stroke. When it persists beyond one month, it is termed persistent; beyond three months for non-traumatic and twelve months for traumatic etiologies, it is termed permanent. Clinical features include spontaneous eye opening with sleep-wake cycles but no purposeful responses. Diagnosis is challenging and may be aided by functional neuroimaging (fMRI) and EEG. Management is supportive; ethical and legal considerations surrounding withdrawal of life-sustaining treatment are complex. Prognosis for recovery of consciousness is poor once permanent criteria are met.

%-------------------------------------------------------------
\section{Epilepsy and Seizure Disorders}
%-------------------------------------------------------------

%-------------------------------------------------------------

%-------------------------------------------------------------

%-------------------------------------------------------------

%-------------------------------------------------------------

%-------------------------------------------------------------

%-------------------------------------------------------------

%-------------------------------------------------------------

\label{sec:Epilepsy and Seizure Disorders}
%-------------------------------------------------------------%-------------------------------------------------------------

\subsection{Epilepsy}
\index{Epilepsy}
\label{sec:Epilepsy}
Epilepsy is a chronic neurological disorder characterized by recurrent, unprovoked seizures resulting from abnormal excessive or synchronous neuronal activity in the brain; it is defined by two or more unprovoked seizures occurring more than 24 hours apart, or one unprovoked seizure with a high recurrence risk. The condition results from genetic, structural (e.g., cortical dysplasia, tumor, stroke), infectious, metabolic, or immune causes. Seizure classification distinguishes focal (partial) from generalized seizures, with further subtypes including absence, tonic-clonic, myoclonic, atonic, and tonic seizures. Clinical features depend on seizure type. Diagnosis relies on clinical history, EEG showing epileptiform discharges, and neuroimaging. Management involves antiseizure medications selected based on seizure type and syndrome classification; drug-resistant epilepsy may be treated with epilepsy surgery, vagus nerve stimulation, or responsive neurostimulation. Prognosis is variable, with many patients achieving seizure control with appropriate therapy.

\subsection{Febrile Seizures}
\index{Febrile Seizures}
\label{sec:Febrile Seizures}
Febrile seizures are the most common seizures in childhood, occurring in 2--5\% of children between 6 months and 5 years of age. The condition results from fever (temperature greater than 38 degrees Celsius) associated with an extracranial infection, in the absence of central nervous system infection or defined electrolyte imbalance. Simple febrile seizures are generalized, last less than 15 minutes, and do not recur within 24 hours; complex febrile seizures are focal, prolonged (greater than 15 minutes), or recurrent. Clinical features include generalized or focal seizure activity triggered by fever. Diagnosis is clinical, after excluding central nervous system infection. Management is supportive; antipyretics prevent recurrence only partially, and prophylactic anticonvulsants are generally not recommended. Prognosis is excellent, with most children outgrowing febrile seizures by age 6.

\subsection{Temporal Lobe Epilepsy}
\index{Temporal Lobe Epilepsy}
\label{sec:Temporal Lobe Epilepsy}
Temporal lobe epilepsy is the most common form of focal epilepsy in adults, often arising from mesial temporal sclerosis (hippocampal sclerosis). The condition results from hippocampal sclerosis, low-grade gliomas, focal cortical dysplasia, or dual pathology. Seizures typically begin with an aura (d\'{e}j\`{a} vu, epigastric rising, fear, olfactory or gustatory hallucinations), followed by impaired awareness, automatisms (lip smacking, hand fumbling), and postictal confusion. Clinical features include focal seizures with impaired awareness. Diagnosis is based on EEG showing temporal interictal epileptiform discharges and MRI revealing hippocampal atrophy and T2/FLAIR signal changes in mesial temporal sclerosis. Management involves antiseizure medications; drug resistance is common, and temporal lobe epilepsy is the most surgically amenable epilepsy, with anterior temporal lobectomy achieving seizure freedom in 60--80\% of selected patients. Prognosis is favorable with surgical treatment for drug-resistant cases.

%-------------------------------------------------------------
\section{Infections of the Brain}
%-------------------------------------------------------------

%-------------------------------------------------------------

%-------------------------------------------------------------

%-------------------------------------------------------------

%-------------------------------------------------------------

%-------------------------------------------------------------

%-------------------------------------------------------------

%-------------------------------------------------------------

\label{sec:Infections of the Brain}
%-------------------------------------------------------------%-------------------------------------------------------------

\subsection{Bacterial Meningitis}
\index{Bacterial Meningitis}
\label{sec:Bacterial Meningitis}
Bacterial meningitis is an acute, life-threatening infection of the leptomeninges and cerebrospinal fluid; the most common causative organisms vary by age: \textit{Streptococcus pneumoniae} and \textit{Neisseria meningitidis} in adults, \textit{Group B Streptococcus} and \textit{Escherichia coli} in neonates, and \textit{Haemophilus influenzae} type b (in unvaccinated populations). The condition results from bacterial invasion of the subarachnoid space. Clinical features include severe headache, high fever, neck stiffness (nuchal rigidity), photophobia, altered mental status, and petechial rash (meningococcal disease); Kernig's and Brudzinski's signs may be positive. Diagnosis requires lumbar puncture showing elevated opening pressure, turbid CSF with neutrophilic pleocytosis, elevated protein, low glucose, and positive Gram stain and culture. Management involves empiric treatment with ceftriaxone plus vancomycin, which should not be delayed; dexamethasone improves outcomes in pneumococcal meningitis. Prognosis depends on the causative organism and timeliness of treatment.

\subsection{Brain Abscess}
\index{Brain Abscess}
\label{sec:Brain Abscess}
A brain abscess is a focal collection of pus within the brain parenchyma, resulting from hematogenous spread, contiguous infection (sinusitis, otitis media, dental infection), or direct inoculation (trauma, neurosurgery); common organisms include \textit{Streptococcus} species, \textit{Staphylococcus aureus}, and anaerobes. The condition results from bacterial infection leading to focal suppuration within brain tissue. Clinical features include headache, fever, focal neurological deficits, seizures, and altered consciousness. Diagnosis is based on MRI with contrast showing a ring-enhancing lesion with surrounding edema and mass effect. Management involves targeted intravenous antibiotics (guided by culture when possible) and, for larger abscesses (greater than 2.5 cm), neurosurgical aspiration or excision. Prognosis has improved significantly with modern imaging and treatment.

\subsection{Neurocysticercosis}
\index{Neurocysticercosis}
\label{sec:Neurocysticercosis}
Neurocysticercosis is the most common parasitic infection of the central nervous system, caused by the larval stage of \textit{Taenia solium} (pork tapeworm); it is endemic in Latin America, Asia, and Africa. The condition results from encystment of larval forms in brain parenchyma, ventricles, or subarachnoid space. Clinical manifestations depend on cyst location and stage and include seizures (the most common presentation worldwide), headache, focal deficits, and hydrocephalus. Diagnosis is supported by neuroimaging showing characteristic cysts (vesicular, colloidal, granular, or calcified stages) and serological testing. Management varies by stage: albendazole and praziquantel for viable cysts with corticosteroid coverage, antiepileptic drugs for seizures, and neurosurgical intervention for hydrocephalus or large symptomatic cysts. Prognosis is generally favorable with appropriate treatment.

\subsection{Viral Encephalitis}
\index{Viral Encephalitis}
\label{sec:Viral Encephalitis}
Viral encephalitis is an inflammation of the brain parenchyma caused by viral infection; herpes simplex virus type 1 (HSV-1) is the most common cause of sporadic fatal encephalitis, typically affecting the temporal and frontal lobes, while other causes include enteroviruses, arboviruses (West Nile, Japanese encephalitis), and varicella-zoster virus. The condition results from direct viral invasion of brain parenchyma. Clinical features include fever, headache, altered consciousness, behavioral changes, seizures, and focal neurological deficits. Diagnosis relies on MRI (temporal lobe involvement in HSV), CSF analysis (lymphocytic pleocytosis, PCR for viral DNA), and EEG. Management involves intravenous acyclovir for HSV encephalitis, which should be initiated empirically. Prognosis is guarded, with significant neurological sequelae in survivors despite treatment.

%-------------------------------------------------------------
\section{Neurodegenerative Diseases}
%-------------------------------------------------------------

%-------------------------------------------------------------

%-------------------------------------------------------------

%-------------------------------------------------------------

%-------------------------------------------------------------

%-------------------------------------------------------------

%-------------------------------------------------------------

%-------------------------------------------------------------

\label{sec:Neurodegenerative Diseases}
%-------------------------------------------------------------%-------------------------------------------------------------

\subsection{Alzheimer's Disease}
\index{Alzheimer's Disease}
\label{sec:Alzheimer's Disease}
Alzheimer's disease is a progressive neurodegenerative disorder and the most common cause of dementia, accounting for 60--70\% of cases. The condition results from the accumulation of amyloid-beta (A$\beta$) plaques between neurons and neurofibrillary tangles composed of hyperphosphorylated tau protein within neurons. Risk factors include advancing age, family history, the apolipoprotein E4 allele, cardiovascular risk factors, and low educational attainment. Clinical features include insidious memory impairment, particularly of recent events, progressing to global cognitive decline, disorientation, language difficulties, and eventual loss of independence. Diagnosis is based on clinical assessment, neuropsychological testing, neuroimaging showing hippocampal atrophy and reduced glucose metabolism on PET, and cerebrospinal fluid biomarkers (A$\beta$42, tau, phospho-tau). Management involves symptomatic treatment with cholinesterase inhibitors (donepezil, rivastigmine, galantamine) and memantine; disease-modifying therapies targeting amyloid remain under investigation. Prognosis is progressive decline with modest benefit from available therapies.

\subsection{Amyotrophic Lateral Sclerosis}
\index{Amyotrophic Lateral Sclerosis}
\label{sec:Amyotrophic Lateral Sclerosis}
Amyotrophic lateral sclerosis is a fatal motor neuron disease characterized by progressive degeneration of both upper and lower motor neurons in the motor cortex, brainstem, and spinal cord. About 10\% of cases are familial, with mutations in SOD1, C9orf72, TARDBP, and FUS genes. The condition results from progressive motor neuron degeneration. Clinical features include focal muscle weakness, atrophy, fasciculations, spasticity, and hyperreflexia; bulbar involvement causes dysarthria and dysphagia. Diagnosis relies on the revised El Escorial criteria, supported by electromyography and exclusion of mimics. Management involves riluzole and edaravone to modestly slow progression, non-invasive ventilation, and multidisciplinary care. Prognosis is poor, with respiratory failure typically causing death within 3--5 years of onset.

\subsection{Dystonia}
\index{Dystonia}
\label{sec:Dystonia}
Dystonia is a movement disorder characterized by sustained or intermittent muscle contractions causing abnormal, often repetitive, movements, postures, or both, with a prevalence of 1 in 3,000. The condition results from dysfunction of the basal ganglia-thalamo-cortical motor circuit, with genetic forms caused by mutations in genes including TOR1A (DYT1, early-onset generalized), THAP1 (DYT6, mixed phenotype), GCH1 (DYT5, dopa-responsive dystonia), and SGCE (DYT11, myoclonus-dystonia). Clinical features vary by distribution: focal (most common, e.g., cervical dystonia/spasmodic torticollis, blepharospasm, oromandibular dystonia, laryngeal dystonia/spasmodic dysphonia, writer's cramp), segmental (involving two contiguous body regions), multifocal, hemidystonia, or generalized. Dystonic movements are typically directional, patterned, and often action-specific, worsened by stress and fatigue, and may be alleviated by sensory tricks (geste antagoniste). Diagnosis is clinical with genetic testing for early-onset or familial cases. Management includes oral medications (trihexyphenidyl, benzodiazepines, baclofen, tetrabenazine), botulinum toxin injections (first-line for focal dystonia), deep brain stimulation of the globus pallidus internus (GPi-DBS, most effective for generalized and cervical dystonia), and levodopa trial for suspected dopa-responsive dystonia. Prognosis: focal dystonia is usually non-progressive after initial spread; generalized dystonia may progress over years but stabilize with treatment.

\subsection{Essential Tremor}
\index{Essential Tremor}
\label{sec:Essential Tremor}
Essential tremor is one of the most common movement disorders, affecting 1\% of the general population and 5\% of individuals over 65 years. The condition results from abnormal oscillatory activity in the cerebello-thalamo-cortical circuit, with genetic heterogeneity (autosomal dominant inheritance in familial cases, ETM1/ETM2 loci). Clinical features include bilateral, symmetric postural and kinetic tremor of the upper limbs (6--12 Hz), which may also involve the head (titubation), voice, chin, or trunk. Tremor is typically absent at rest, worsens with goal-directed movement, and is exacerbated by stress, fatigue, and caffeine. Alcohol consumption temporarily reduces tremor amplitude in 50--70\% of patients. Diagnosis is clinical, based on the Movement Disorder Society consensus criteria and exclusion of other causes (Parkinson disease, dystonic tremor, enhanced physiologic tremor, drug-induced tremor). Management includes first-line treatment with propranolol or primidone, alternative agents (topiramate, gabapentin, benzodiazepines), and for medication-refractory cases, deep brain stimulation of the ventral intermediate nucleus (VIM) of the thalamus or focused ultrasound thalamotomy. Prognosis is benign but slowly progressive; disability from tremor may be significant.

\subsection{Friedreich's Ataxia}
\index{Friedreich's Ataxia}
\label{sec:Friedreich's Ataxia}
Friedreich's ataxia is the most common hereditary ataxia, caused by homozygous GAA trinucleotide repeat expansion in the FXN gene encoding frataxin, with a prevalence of 1 in 50,000. The condition results from frataxin deficiency, leading to mitochondrial iron accumulation, oxidative stress, and impaired energy metabolism, primarily affecting the dorsal root ganglia, posterior columns, corticospinal tracts, and cerebellum. Clinical features include progressive gait and limb ataxia, dysarthria, areflexia, loss of vibratory and proprioceptive sensation, extensor plantar responses, scoliosis, pes cavus, hypertrophic cardiomyopathy, and increased risk of diabetes mellitus. Onset is typically in childhood or adolescence (mean age 10--15 years). Diagnosis is based on clinical criteria and confirmed by FXN genetic testing. Management involves symptomatic treatment: physical therapy for balance and mobility, occupational therapy, speech therapy, and orthopedic management of scoliosis. Cardiac monitoring is essential (echocardiography, ECG). Prognosis is progressive; most patients lose ambulation 10--15 years after symptom onset, with mean survival to age 40--50, primarily due to cardiomyopathy.

\subsection{Frontotemporal Dementia}
\index{Frontotemporal Dementia}
\label{sec:Frontotemporal Dementia}
Frontotemporal dementia encompasses a group of neurodegenerative disorders selectively affecting the frontal and temporal lobes; onset typically occurs between ages 40 and 65, making it a common cause of early-onset dementia. The condition results from accumulation of tau, TDP-43, or FUS proteins. Clinical features include behavioral variant FTD with personality changes, disinhibition, apathy, and executive dysfunction; language variants include semantic dementia (progressive loss of word meaning) and non-fluent or agrammatic primary progressive aphasia. Diagnosis is based on MRI showing disproportionate frontal and temporal lobe atrophy. Management is supportive, focusing on behavioral management, speech therapy, and caregiver support. Prognosis is progressive, with no disease-modifying treatment available.

\subsection{Huntington's Disease}
\index{Huntington's Disease}
\label{sec:Huntington's Disease}
Huntington's disease is an autosomal dominant neurodegenerative disorder caused by a CAG trinucleotide repeat expansion in the huntingtin (HTT) gene on chromosome 4. The condition results from mutant huntingtin protein causing progressive degeneration of neurons, particularly in the caudate nucleus and putamen (striatum). Onset typically occurs between ages 30 and 50, with disease severity correlating inversely with the number of CAG repeats. Clinical features include a triad of motor dysfunction (chorea, dystonia, bradykinesia), cognitive decline (executive dysfunction, dementia), and psychiatric disturbances (depression, irritability, psychosis). Diagnosis is confirmed by genetic testing. Management involves symptomatic treatment with tetrabenazine or deutetrabenazine for chorea, antipsychotics, and antidepressants; genetic counseling is essential. Prognosis is progressive and fatal, with no curative treatment available.

\subsection{Motor Neuron Disease}
\index{Motor Neuron Disease}
\label{sec:Motor Neuron Disease}
Motor neuron disease (MND) is a group of progressive neurodegenerative disorders affecting upper motor neurons (UMNs), lower motor neurons (LMNs), or both, with a lifetime risk of 1 in 300. The spectrum includes amyotrophic lateral sclerosis (ALS, most common, both UMN and LMN), progressive bulbar palsy (PBP, predominantly bulbar onset), progressive muscular atrophy (PMA, predominantly LMN), and primary lateral sclerosis (PLS, predominantly UMN). The condition results from TDP-43 pathology, glutamate excitotoxicity, oxidative stress, mitochondrial dysfunction, and impaired protein degradation, with genetic causes including C9orf72 hexanucleotide repeat expansion, SOD1, TARDBP, FUS, and TBK1 mutations. Clinical features include progressive muscle weakness, wasting, fasciculations, spasticity, hyperreflexia (UMN signs), and bulbar symptoms (dysarthria, dysphagia). ALS presentations include limb-onset (70\%), bulbar-onset (25\%), and respiratory-onset (5\%). Cognitive impairment is present in 30--50\% (frontotemporal dementia in 15\%). Diagnosis is based on revised El Escorial criteria, electromyography showing widespread denervation, and exclusion of mimics. Management involves multidisciplinary care, riluzole (extends survival by 2--3 months), edaravone (slows functional decline in early ALS), noninvasive ventilation for respiratory insufficiency, and symptomatic treatments (sialorrhea, spasticity, pain). Prognosis is uniformly fatal in ALS: median survival 3--5 years from symptom onset, with respiratory failure as the most common cause of death.

\subsection{Multiple Sclerosis}
\index{Multiple Sclerosis}
\label{sec:Multiple Sclerosis}
Multiple sclerosis is a chronic autoimmune demyelinating disease of the central nervous system characterized by inflammatory demyelination and axonal injury occurring at different times and locations; the relapsing-remitting form (RRMS) is most common, while progressive forms (primary progressive, secondary progressive) involve gradual neurological deterioration. The condition results from autoimmune-mediated demyelination. Clinical features include episodic neurological deficits such as optic neuritis, transverse myelitis, brainstem syndromes, and sensory or motor disturbances. Diagnosis is based on MRI revealing periventricular, juxtacortical, infratentorial, and spinal cord T2-hyperintense lesions, often with gadolinium enhancement in active plaques; cerebrospinal fluid may show oligoclonal bands. Management involves disease-modifying therapies including interferons, glatiramer acetate, natalizumab, fingolimod, ocrelizumab, and others; acute relapses are treated with high-dose corticosteroids or plasmapheresis. Prognosis varies widely depending on disease course and treatment response.

\subsection{Parkinson's Disease}
\index{Parkinson's Disease}
\label{sec:Parkinson's Disease}
Parkinson's disease is a chronic, progressive neurodegenerative disorder affecting the dopaminergic neurons of the substantia nigra pars compacta in the midbrain. The condition results from loss of pigmented dopaminergic neurons and the presence of Lewy bodies composed of misfolded alpha-synuclein. Clinical features include cardinal motor manifestations---bradykinesia, resting tremor, rigidity, and postural instability---as well as non-motor symptoms such as anosmia, constipation, REM sleep behavior disorder, depression, and autonomic dysfunction, which often precede motor manifestations by years. Diagnosis is clinical, supported by response to dopaminergic therapy and DaT-SPECT imaging. Management involves dopamine replacement with levodopa/carbidopa, dopamine agonists, MAO-B inhibitors, and COMT inhibitors; advanced disease may be managed with deep brain stimulation of the subthalamic nucleus or globus pallidus interna. Prognosis is progressive, with disease-modifying therapies remaining under investigation.

%-------------------------------------------------------------
\section{Peripheral Nerve Diseases}
%-------------------------------------------------------------

%-------------------------------------------------------------

%-------------------------------------------------------------

%-------------------------------------------------------------

%-------------------------------------------------------------

%-------------------------------------------------------------

%-------------------------------------------------------------

%-------------------------------------------------------------

\label{sec:Peripheral Nerve Diseases}
%-------------------------------------------------------------%-------------------------------------------------------------

\subsection{Bell's Palsy}
\index{Bell's Palsy}
\label{sec:Bell's Palsy}
Bell's palsy is an acute, idiopathic, unilateral facial nerve (cranial nerve VII) paralysis of lower motor neuron type and is the most common cause of acute facial palsy. The condition is thought to result from reactivation of herpes simplex virus type 1. Onset is typically abrupt, with complete facial weakness developing over hours to days. Clinical features include inability to close the affected eye, drooping mouth, hyperacusis, and loss of taste on the anterior two-thirds of the tongue. Diagnosis is clinical, after excluding causes such as Ramsay Hunt syndrome, Lyme disease, and stroke. Management with oral corticosteroids (prednisone) started within 72 hours of onset significantly improves outcomes; eye protection is essential to prevent corneal damage. Prognosis is favorable, with most patients recovering fully within weeks to months.

\subsection{Carpal Tunnel Syndrome}
\index{Carpal Tunnel Syndrome}
\label{sec:Carpal Tunnel Syndrome}
Carpal tunnel syndrome is the most common entrapment neuropathy, caused by compression of the median nerve as it passes through the carpal tunnel in the wrist. Risk factors include repetitive wrist movements, pregnancy, hypothyroidism, diabetes, rheumatoid arthritis, and space-occupying lesions. The condition results from mechanical compression of the median nerve within the carpal tunnel. Clinical features include numbness and tingling in the median nerve distribution (thumb, index, middle, and radial half of the ring finger), nocturnal pain, and weakness of the thenar muscles in advanced cases. Diagnosis is based on clinical examination (Phalen's test, Tinel's sign) and nerve conduction studies confirming delayed median nerve conduction across the wrist. Management begins with wrist splinting (especially at night), ergonomic modifications, and corticosteroid injection; surgical carpal tunnel release is indicated for severe or refractory cases. Prognosis is excellent with appropriate treatment.

\subsection{Charcot-Marie-Tooth Disease}
\index{Charcot-Marie-Tooth Disease}
\label{sec:Charcot-Marie-Tooth Disease}
Charcot-Marie-Tooth disease (CMT) is the most common hereditary peripheral neuropathy, affecting 1 in 2,500 individuals. It comprises a genetically heterogeneous group of disorders classified by nerve conduction velocity and inheritance pattern: CMT1 (demyelinating, NCV <38 m/s, autosomal dominant, most commonly PMP22 duplication), CMT2 (axonal, NCV >38 m/s, autosomal dominant, MFN2 most common), CMTX (X-linked, GJB1 mutations), and CMT4 (autosomal recessive). Clinical features include slowly progressive distal muscle weakness and atrophy (beginning in peroneal muscles, leading to foot drop and steppage gait), pes cavus, hammer toes, distal sensory loss, and diminished or absent deep tendon reflexes. Upper extremity involvement occurs later, producing intrinsic hand muscle weakness and claw hand deformity. Diagnosis is based on clinical examination, nerve conduction studies, and genetic testing (next-generation sequencing panels). Management is supportive: physical and occupational therapy, ankle-foot orthoses for foot drop, surgical correction of foot deformities, and pain management for neuropathic pain. Prognosis is generally good; most patients remain ambulatory with normal life expectancy, though the severity varies widely by subtype.

%-------------------------------------------------------------
\section{Traumatic Brain Injury}
%-------------------------------------------------------------

%-------------------------------------------------------------

%-------------------------------------------------------------

%-------------------------------------------------------------

%-------------------------------------------------------------

%-------------------------------------------------------------

%-------------------------------------------------------------

%-------------------------------------------------------------

\label{sec:Traumatic Brain Injury}
%-------------------------------------------------------------%-------------------------------------------------------------

\subsection{Chronic Traumatic Encephalopathy}
\index{Chronic Traumatic Encephalopathy}
\label{sec:Chronic Traumatic Encephalopathy}
Chronic traumatic encephalopathy is a progressive neurodegenerative disease associated with repetitive head impacts, found predominantly in contact sport athletes, military personnel exposed to blast injuries, and victims of repeated physical abuse. The condition results from accumulation of hyperphosphorylated tau protein in neurons and astrocytes, typically in a perivascular distribution at the depths of cortical sulci. Clinical features include behavioral and mood changes (depression, apathy, irritability, impulsivity), cognitive impairment (progressive dementia), and motor abnormalities (parkinsonism); symptoms typically manifest years to decades after the repetitive head trauma. Diagnosis is currently only confirmed postmortem by neuropathological examination. Management is supportive, as no disease-modifying treatment exists. Prognosis is progressive and incurable; the discovery of chronic traumatic encephalopathy has led to significant changes in sports safety protocols.

\subsection{Concussion}
\index{Concussion}
\label{sec:Concussion}
Concussion is a mild traumatic brain injury caused by a direct blow to the head or a biomechanical force transmitted to the brain, resulting in a transient disturbance of neurological function; it is the most common form of traumatic brain injury. The condition results from a complex neurometabolic cascade involving ionic flux, neurotransmitter release, impaired cerebral blood flow regulation, and axonal injury, without macrostructural damage visible on conventional neuroimaging. Clinical features include headache, dizziness, confusion, amnesia, visual disturbances, sensitivity to light and noise, difficulty concentrating, and emotional lability; loss of consciousness occurs in fewer than 10\% of cases. Diagnosis is clinical; CT head is indicated to exclude intracranial hemorrhage in patients with risk factors (age greater than 65, anticoagulant use, focal neurological signs, persistent vomiting, worsening headache). Management involves physical and cognitive rest for 24--48 hours, followed by a graduated return-to-activity protocol. Prognosis is favorable, although post-concussion syndrome (persistent symptoms beyond 4 weeks) occurs in 10--30\% of patients; second impact syndrome is a rare but potentially fatal complication.

\subsection{Epidural Hematoma}
\index{Epidural Hematoma}
\label{sec:Epidural Hematoma}

An epidural hematoma is a traumatic intracranial hemorrhage accumulating between the inner skull table and the dura mater, most commonly resulting from skull fracture crossing the pterion with laceration of the middle meningeal artery. Clinical features include transient loss of consciousness at injury followed by a lucid interval (lasting hours), after which neurological deterioration ensues with severe headache, contralateral hemiparesis, ipsilateral dilated and unreactive pupil (due to uncal herniation with CN III compression), and Cushing triad (hypertension, bradycardia, irregular respiration). Diagnosis is by non-contrast head CT showing a convex, biconvex (lens-shaped) hyperdense collection that does not cross cranial suture lines. Management requires urgent neurosurgical craniotomy and hematoma evacuation to prevent fatal brain herniation. Prognosis is excellent if evacuated promptly before herniation.

\subsection{Subdural Hematoma}
\index{Subdural Hematoma}
\label{sec:Subdural Hematoma}
A subdural hematoma is a traumatic intracranial bleeding accumulating in the potential space between the dura mater and arachnoid mater, resulting from tearing of bridging veins. Acute subdural hematomas follow high-impact trauma; chronic subdural hematomas occur predominantly in elderly patients and alcoholics, where cerebral atrophy stretches bridging veins, making them susceptible to rupture from minor head trauma. Clinical features include headache, confusion, fluctuating consciousness, and focal neurological deficits, developing over hours in acute cases or weeks/months in chronic cases. Diagnosis is established by non-contrast head CT: acute SDH appears crescent-shaped hyperdense; chronic SDH appears hypodense; SDH crosses suture lines but does not cross dural reflections (falx/tentorium). Management of large or symptomatic hematomas involves surgical evacuation (burr holes or craniotomy); small asymptomatic SDHs are managed conservatively. Prognosis depends on underlying brain injury.

%-------------------------------------------------------------
\section{Vestibular Disorders}
%-------------------------------------------------------------

%-------------------------------------------------------------

%-------------------------------------------------------------

%-------------------------------------------------------------

%-------------------------------------------------------------

%-------------------------------------------------------------

%-------------------------------------------------------------

%-------------------------------------------------------------

\label{sec:Vestibular Disorders}
%-------------------------------------------------------------%-------------------------------------------------------------

\subsection{Vertigo}
\index{Vertigo}
\label{sec:Vertigo}
Vertigo is the sensation of spinning or movement of the self or surroundings when there is no actual movement; it is a symptom, not a disease, and may originate from peripheral (inner ear) or central (brainstem, cerebellar) causes. Peripheral vertigo is more common and results from dysfunction of the vestibular apparatus or nerve; causes include benign paroxysmal positional vertigo, vestibular neuritis, labyrinthitis, and M\'{e}ni\`{e}re's disease. Central vertigo results from brainstem or cerebellar pathology; causes include posterior circulation stroke, vestibular migraine, multiple sclerosis, and cerebellar tumors. Clinical features of peripheral vertigo include severe rotational vertigo with nystagmus (beating away from the affected side), nausea, and vomiting; central vertigo tends to be less severe but may be continuous with diplopia, dysarthria, and limb ataxia. Diagnosis involves detailed history, neurological examination, Dix-Hallpike test for BPPV, and neuroimaging when central causes are suspected. Management depends on the underlying cause: canalith repositioning maneuvers for BPPV, vestibular rehabilitation for vestibular neuritis, and stroke protocols for central vertigo; vestibular suppressants provide symptomatic relief but should be used short-term. Prognosis depends on the underlying cause.

